\subsection{Building intuition with four entities}
Imagine four entities able to extract a resource from their
environment. Each has a mobilizable stock $K_i$: the larger it is,
the more it extracts, but each additional unit of stock yields less
than the previous one. The function $f_i(K)=A_iK^{\gamma_i}$
describes this \emph{extraction technology}. The input joins the
stock, a fraction of which decays at each step. New entities appear
with the same endowment $K^\circ$ and, in the reference regime,
the same production exponent $\gamma^\circ=1/2$.

A well-endowed entity can lend stock to another, where that stock
produces more at the margin. This real transfer simultaneously
creates a claim for the lender and a debt for the borrower. The
\emph{principal} $q$ is the nominal amount recorded; the \emph{interest
service} $rq$ is the payment due at each step. Interest moves the
resource between entities, without creating any. The principal is not
progressively repaid in this model: this is a strong institutional
assumption.

\begin{figure}[htbp]\centering
\begin{tikzpicture}[>=Stealth,font=\small,
 ent/.style={circle,draw=blue!60!black,fill=blue!7,minimum size=15mm,align=center},
 lab/.style={font=\scriptsize,fill=white,inner sep=2pt}]
\draw[rounded corners,fill=gray!3,draw=gray!50] (-1.5,-1.7) rectangle (9.5,2.3);
\node[ent] (a) at (0,0.7) {$e_1$\\$K_1=16$};
\node[ent] (b) at (4,0.7) {$e_2$\\$K_2=4$};
\node[ent] (c) at (7.5,1.2) {$e_3$\\$K_3=9$};
\node[ent] (d) at (7.5,-0.8) {$e_4$\\$K_4=1$};
\draw[->,very thick,blue!70!black] (a) to[bend left=18]
 node[lab,above] {initial loan $q=6$} (b);
\draw[->,dashed,thick,orange!80!black] (b) to[bend left=22]
 node[lab,below] {interest $rq$, at each step} (a);
\draw[->,green!45!black,thick] (-.9,2.8)--(a);
\node[lab] at (2.4,2.8) {input extracted from the environment: $f_i(K_i)$};
\draw[->,green!45!black,thick] (7.5,2.8)--(c);
\draw[->,red!70!black] (d)--(9.8,-1.4) node[right,font=\scriptsize] {erosion};
\node[align=center,font=\scriptsize] at (2.7,-1.25)
 {After the loan: $K_1=K_2=10$; $C_1=6$, $D_2=6$.\\
 The net worths remain $NW_1=16$ and $NW_2=4$.};
\end{tikzpicture}
\caption{\textbf{Real stock and accounting commitment are not the same thing.}
Fictitious example, inspired by the pedagogical diagram of the internship
presentation. The initial values are chosen to explain the mechanism, not
extracted from a simulation. Entities 3 and 4 show that not all entities
have to take part in the same exchange. Inputs and erosion are drawn on
only a few entities to lighten the diagram; the rules apply to all.
The external reservoir is not simulated.}
\label{fig:exemple}
\end{figure}

The loan thus preserves the instantaneous net worth of each partner,
but changes its productive capacity and its future commitments. An
entity may later be unable to honor its service, or owe more than it
holds. Its disappearance wipes out the claims held on it: its lenders
can in turn become insolvent. The coupling between real stocks,
persistent commitments and propagated losses is the mechanism studied
here.
